\section{A diffuse model with directly executable paired certificates}\label{sec:reset50}
The general transfer theorem allows controlled densities. A reset subclass permits a direct integrated lower bound and avoids paying a conservative kernel-approximation error. Let $X=[0,1]$, let $\nu$ and every transition kernel be uniform on $[0,1]$, and take
\[
 r_t(x,a)=\gamma_a g_t(x),\qquad 0\le\gamma_a\le1,
 \quad\max_a\gamma_a=1,
\]
with positive affine $g_t$ and nonnegative affine implementation costs $k_{ta}$. The terminal reward is zero. Transitions have a continuous density and a diffuse initial law. They are independent of state and action, so this subclass exploits a reset structure rather than test the generic controlled-kernel search.

The operating advantage is $d_{ta}(x)=(1-\gamma_a)g_t(x)$. Define cumulative regret moments
\begin{equation}\label{eq:reset-moments50}
 M_T=0,\qquad
 M_t=\int_0^1\sum_a p_t(a\mid x)d_{ta}(x)\,dx+\beta M_{t+1}.
\end{equation}
The continuation term in this equation is essential; $M_t$ is not merely the one-period expected disadvantage. The original contract and objective become
\begin{align}
 \sum_a p_t(a\mid x)d_{ta}(x)+\beta M_{t+1}&\le\varepsilon
 &&\text{for all }x,t,\label{eq:reset-contract50}\\
 C_0^p(\nu)&=\sum_{t=0}^{T-1}\beta^t\int_0^1\sum_a p_t(a\mid x)k_{ta}(x)\,dx.\label{eq:reset-objective50}
\end{align}
Equations~\eqref{eq:reset-moments50}--\eqref{eq:reset-objective50} are linear in the measurable probabilities and moments. They retain one common continuation moment at every date, rather than an independent choice of that moment for each restart.

\begin{proposition}[Integrated lower certificate for all Borel policies]\label{prop:reset-dual50}
Choose arbitrary nonnegative bounded Borel multipliers $\lambda_t(x)$ and arbitrary real numbers $q_t$. Set $\Lambda_t=\int\lambda_t(x)dx$ and
\[
 s_0=-q_0,\qquad s_t=-q_t+\beta(q_{t-1}+\Lambda_{t-1})\quad (t\ge1).
\]
Then the following is a lower bound on the original unrestricted Borel-policy optimum:
\begin{equation}\label{eq:reset-dual50}
 \begin{split}\mathcal L(\lambda,q)=\max\bigg\{0,
 &\sum_t\int_0^1\min_a\{\beta^tk_{ta}(x)+(\lambda_t(x)+q_t)d_{ta}(x)\}\,dx\\
 &-\varepsilon\sum_t\Lambda_t+\varepsilon\sum_t\min(0,s_t)\bigg\}.\end{split}
\end{equation}
For stepwise constant $\lambda_t$ and rational affine primitives, its integrals can be evaluated with finite rational operations.
\end{proposition}
\begin{proof}
For a feasible policy, multiply \eqref{eq:reset-contract50} minus $\varepsilon$ by $\lambda_t(x)$ and integrate. Add $q_t$ times \eqref{eq:reset-moments50} written as $\int p_td_t+\beta M_{t+1}-M_t=0$. These additions do not increase the feasible objective. The coefficient of $M_t$ is $s_t$. Every feasible moment lies in $[0,\varepsilon]$, so its contribution is at least $\varepsilon\min(0,s_t)$. Minimizing each policy row over the simplex gives the pointwise minimum in \eqref{eq:reset-dual50}. This inequality holds for every feasible Borel policy before any cell restriction. Nonnegative costs permit clipping the lower bound at zero.

On a cell where $\lambda_t$ is constant, each expression inside the minimum is affine. Partition the cell further at pairwise line intersections; one line minimizes on each resulting open interval. Its integral is the rational trapezoidal integral. Endpoint ties have zero Lebesgue measure and do not change the value.\end{proof}

\subsection{Constructing a globally feasible upper endpoint}
Partition $[0,1]$ into $N$ equal intervals and use one policy row $p_{tca}$ on each cell. Let $\bar d_{tca}$ and $\bar k_{tca}$ be their exact cell averages, and let $d^{\max}_{tca}$ be the maximum of the two endpoint disadvantages. The following finite LP generates a proposal:
\begin{align}
 \min_{p,M}\quad &\sum_{t,c,a}\frac{\beta^t}{N}p_{tca}\bar k_{tca},\nonumber\\
 M_t&=N^{-1}\sum_{c,a}p_{tca}\bar d_{tca}+\beta M_{t+1},\nonumber\\
 \sum_a p_{tca}d^{\max}_{tca}+\beta M_{t+1}&\le\varepsilon,
 \quad p_{tc}\in\Delta_m,\quad M_T=0.\label{eq:reset-lp50}
\end{align}
The sums of actionwise maxima can be conservative when slopes differ. They guarantee the pointwise constraint without assuming a cellwise optimum of the unrestricted problem. The LP has $TNm+T$ variables and $2TN+T$ stated equality and inequality rows, before simple variable bounds.

Numerical probabilities are reconstructed as dyadic rationals and normalized. Process dates backward using the actual moment of the already reconstructed successor. The remaining allowance is $a_t=\varepsilon-\beta M_{t+1}>0$. If a row's expected endpoint-maximum disadvantage exceeds $a_t$, mix toward an action with $\gamma_a=1$ by the exact amount needed. This action has zero disadvantage. Round all other probabilities downward to a dyadic grid and transfer the remainder to the zero-disadvantage action. The row remains feasible. Recompute the exact integrated moment using \eqref{eq:reset-moments50}, not the one-period integral alone. Since every row satisfies \eqref{eq:reset-contract50}, integrating proves $M_t\le\varepsilon$ and justifies the next step.

The resulting piecewise constant Borel policy provides an exactly integrated upper cost. The lower certificate uses the same continuous primitives but is not restricted to this policy class. A numerical LP status is neither necessary nor sufficient for accepting an endpoint: a failed numerical solve can supply a weaker proposal, and only the reconstructed inequalities determine validity.

\subsection{Arithmetic, stopping, and approximation}
Policy and multiplier proposals use 40-bit dyadic reconstruction. Each affine-envelope integral is evaluated rationally and rounded \emph{downward} to 50 binary fractional bits before accumulation. This adds less than $TN2^{-50}$ of total downward error, prevents unnecessary denominator growth, and leaves the lower inequality valid. The upper endpoint is a rational integral of the deployed policy, without downward rounding. The independent readers repeat these conventions.

The operational stopping condition is the original continuous-model inequality $U-\mathcal L\le10^{-3}$. The same units apply to every primary model. Relative width and width normalized by the cost of the always-preferred action are secondary reports and do not change target status. A lower bound on implementation savings relative to that feasible reference is its exact cost minus $U$. A fixed administration fee is worth incurring only when it is smaller than a justified cost saving against the specified comparator; no dollar conversion is built into the certificate.

Refining the cells can be justified independently of a numerical solver's dual selection. Here $a_P=0$ in Theorem~\ref{thm:density-pair50}; affine reward and cost approximation errors are explicitly of order $N^{-1}$. The cell model has action-independent transitions, so its common-policy moment formulation is a finite LP. Solving the two corrected-allowance finite programs and using Corollary~\ref{cor:density-stop50} gives a complete paired approximation algorithm for this subclass. The direct integrated dual often gives a substantially narrower interval than that sufficient transfer bound. We accept the direct interval only through its own certificate; we do not assume that every sequence of floating LP marginals must converge.
